A eq. (3.24) é dimensionalmente consistente quando cada termo da equação compartilha exatamente as mesmas dimensões físicas fundamentais.

\textbf{Condições necessárias e suficientes:}
\begin{enumerate}
\item \textbf{Homogeneidade dimensional:} Todos os termos somados ou subtraídos devem ter as mesmas dimensões. Se a equação é do tipo $A + B = C$, então $[A] = [B] = [C]$.

\item \textbf{Consistência dos parâmetros:} Quaisquer coeficientes numéricos (adimensionais) podem aparecer livremente, mas parâmetros físicos que acompanham os termos devem contribuir para as dimensões corretas.

\item \textbf{Verificação explícita para tensão mecânica:} Para equações envolvendo tensão $\sigma$ (força por área), a consistência dimensional exige que:
\[\left[\sigma\right] = \frac{F}{A} = \frac{M\,L\,T^{-2}}{L^2} = M\,L^{-1}\,T^{-2}.\]
Todo termo que representa tensão deve possuir estas dimensões.
\end{enumerate}

\textbf{Resumo:} A eq. (3.24) é dimensionalmente consistente precisamente quando as grandezas substituídas satisfazem a condição de que os expoentes de $M$, $L$ e $T$ são idênticos em todos os seus termos.

